% 按真实问题数写作；本文件仅为不可直接交付的槽位模板。
% 先读当前gmcm-abstract与摘要题型规则，建立abstract_evidence_path。
% 不反向造结果；fragment不写标题/关键词，关键词只由manifest生成。
本文针对【核心科学任务】，建立【总体模型】，采用\textbf{【核心方法】}，得到\textbf{【真实主要结论】}及\textbf{【关键数值与单位；解析结果不强造数字】}。

\textbf{针对问题一}，【描述任务；仅优化题强调目标与约束】，采用\textbf{【核心方法名称】}，得到\textbf{【主要结论短语】}，关键结果为\textbf{【真实数值及单位】}。【将本问独立验证或解析证明及关键结果融入此段】。

% 按真实Q继续；只有真实依赖前问才写依赖，不机械复制三个问题。
% 方法差异/创新证据可融入对应问题或用1--2句单独说明；无证据不自夸。
% 不写“最后...”局限/未来工作/内部验收/交付结尾。
